\section{Results and Current Evidence}

\subsection{Watermarked PushT}

The strongest completed result is a sustained geometric reversal in the
watermarked PushT setting.  At the end of the full seed-zero run, zero-floor cosine
\shortmethod{} yields same-content cosine $0.807$ and same-tag cosine $0.176$,
for a margin of $+0.631$.  Its final planning success is $0.80$, the mean of the
last ten evaluations is $0.774$, and the best 50-episode checkpoint reaches
$0.88$.  The peak is not used as the primary estimate.

Frozen probes show why this is not simply information deletion.  Relative to
a prediction-weight control, the projection-space block-angle probe remains
essentially unchanged ($0.800$ versus $0.795$), while tag RGB $R^2$ decreases
from $0.889$ to $0.502$.  The pusher and block positions also remain strongly
decodable.  The intervention therefore changes relative geometric prominence
more clearly than raw task-variable availability.

At convergence, the mean prediction/control cosine is approximately $0.072$
and only about $0.106$ of the raw prediction-gradient norm is admitted.  This
explains both the watermark resistance and the central concern of this paper:
cosine admission may act as an aggressive encoder-side prediction bottleneck.

\subsection{Magnitude versus direction}

In matched 10k screening runs, norm-matched scalar attenuation reaches success
$0.42$, cosine routing $0.36$, and a retention-matched shuffled guide $0.28$.
The scalar result shows that a substantial fraction of the early benefit can
come from limiting prediction-gradient magnitude.  The shuffled-guide gap is
preliminary evidence that correct sample-specific correspondence also matters.
Full schedules and multiple seeds are required before assigning either effect
priority.

\subsection{Cross-task snapshot}

Table~\ref{tab:cross-task} records completed seed-zero checkpoints.  The rows
mix 10k screening and full schedules and therefore document current evidence,
not a final leaderboard.

\begin{table}[t]
\centering
\small
\caption{Current matched seed-zero planning success. ``Full'' and ``10k'' are
not compared across rows. Reacher uses the recovered historical evaluation
protocol at epoch 10.}
\label{tab:cross-task}
\begin{tabular}{llccc}
\toprule
Task & Condition & Schedule & LeWM & Cosine \shortmethod{} \\
\midrule
TwoRooms & clean  & 10k  & 0.88 & 0.90 \\
TwoRooms & tagged & 10k  & 0.60 & 0.98 \\
Cube     & clean  & full & 0.70 & 0.76 \\
Cube     & tagged & full & 0.74 & 0.84 \\
Reacher  & clean  & epoch 10 & 0.90 & 0.86 \\
\bottomrule
\end{tabular}
\end{table}

TwoRooms and Cube do not reproduce the same catastrophic tagged-baseline
collapse as PushT.  They remain useful as robustness checks, but cannot by
themselves support a universal nuisance-removal claim.  Reacher exposes the
opposite failure boundary: aggressive cosine routing can suppress predictive learning that
the planner still needs.

\subsection{Reacher and residual admission}

A one-epoch diagnostic under the historical evaluator gives success $0.78$
for LeWM, $0.70$ for cosine routing, and $0.82$ for both control-only and
$\beta=0.75$ soft admission.  Each number is based on only 50 evaluation
episodes and is therefore a screening signal, not a final result.  The full
$\beta=0.75$ Reacher run and the paired tagged-PushT guardrail are in progress.
Their joint outcome decides whether a residual prediction path repairs the
failure or simply reintroduces the predictable shortcut.

\subsection{Static but decision-causal goals}

RandGoal shows that within-episode inverse dynamics is an incomplete guide for
static task variables.  A goal can change the expert action distribution while
remaining absent from short local transition residuals.  Current action-only
routing does not yet establish a reliable improvement in this setting.  We
retain RandGoal as a boundary condition rather than hiding it: suppressing all
static information would solve the watermark benchmark for the wrong reason.

\subsection{Pending mechanism tests}

The conflict-only PushT run tests whether the relatively rare negative
gradient conflicts explain the gain.  The component-gradient diagnostic
measures individual horizons and their joint span on tagged PushT and clean
Reacher.  We will replace this paragraph with complete curves and confidence
intervals after the predeclared runs finish.  \todo{Insert conflict-only,
soft75 guardrail, and component-span results; then promote selected variants
to seeds 1--2.}
